\section{Related Work}
\label{sec:related_work}

The literature on parameter-efficient multi-task ensembling can be broadly categorized into static weight merging, parametric dynamic routing, and non-parametric activation blending \cite{yang2024modelmerging}.

\subsection{Static Weight Merging}
To combine specialized models without increasing test-time parameters, static weight merging has emerged as a popular paradigm. Early methods such as Uniform Weight Merging \cite{wortsman2022modelmerging} simply average the parameters of multiple task-specific expert networks. To prevent representational interference and parameter conflict, more sophisticated techniques have been developed. For example, Task Arithmetic \cite{ilharco2022taskarithmetic} treats fine-tuned weights as task vectors that can be added or subtracted. TIES-Merging \cite{yadav2023ties} resolves sign conflicts and prunes low-magnitude parameter changes, while DARE \cite{yu2023dare} employs randomized delta-weight pruning with scaling to preserve base performance. While computationally elegant, these static methods are fundamentally limited by representation collapse: when experts are trained on highly heterogeneous or disjoint task domains, statically averaging their parameters compromises task-specific performance, as we demonstrate in our uniform merging baseline.

\subsection{Parametric Dynamic Routing \& Mixture-of-Experts}
Mixture-of-Experts (MoE) architectures \cite{shazeer2017moe, fedus2022switch} introduce dynamic ensembling by training parametric gating networks (routers) to route tokens to specialized networks on-the-fly. While effective, standard MoE models must be trained from scratch on massive datasets and introduce significant parameter bloat. 
To enable dynamic ensembling over parameter-efficient specialized adapters (such as LoRAs), several parametric routers have been proposed. Methods like Linear Router \cite{pfsr2024} train a dedicated linear gating layer on low-data calibration splits ($B_{\text{cal}} = 64$) to output task-specialization coefficients. However, these parametric routers overfit aggressively to the low-data calibration manifold. Under true batch-independent streaming regimes ($B=1$), they experience \emph{Vectorization Collapse} because their trained decision boundaries fail to generalize to individual samples without batch-average normalization.

\subsection{Non-Parametric Activation Blending}
To overcome the overfitting and training overhead of parametric routers, recent research has focused on non-parametric activation ensembling. These frameworks combine specialized expert representations directly in the activation space during the forward pass.
\begin{itemize}
    \item \textbf{PFSR with Micro-Batch Homogenization (MBH):} PFSR \cite{pfsr2024} performs non-parametric routing using the base model's final classification heads. However, to handle heterogeneous streams, it relies on MBH scheduling, which groups incoming queries into homogeneous micro-batches. This introduces high systems complexity and $O(K)$ sequential backbone passes, making it impractical for latency-critical edge devices.
    \item \textbf{SPS-ZCA:} Zero-Shot Centroid Alignment (ZCA) \cite{spszca2025} aligns task activations at specific layers. However, to prevent representational drift during routing, it restricts expert adapters to mid-to-late layers, creating a training-test mismatch since experts were originally trained across all layers.
    \item \textbf{SABLE:} SABLE \cite{sable2025} performs sample-wise activation ensembling by computing routing coefficients in later layers of the network. To resolve the Routing Paradox, it employs Late Adaptation, freezing and leaving the first 10 layers of the network completely unadapted. While SABLE successfully mitigates representation collapse compared to static uniform merging, leaving early-to-mid layers unadapted prevents it from utilizing specialized early-layer features, limiting its representation capacity.
\end{itemize}

Unlike these prior ensembling frameworks, PEAR relentlessly applies Occam's razor. It resolves the Routing Paradox and the Early-Feature Loss Trade-Off simultaneously by performing dynamic ensembling inside the frozen Patch Embedding layer (Layer 0). By using Layer 0 as a zero-cost, non-parametric feature extractor, PEAR computes highly robust routing coefficients before propagating through the network, allowing expert LoRA adapters to remain fully active across 100\% of the layers with flat $O(1)$ latency.
